\documentclass[10pt,a4paper]{amsart}
\usepackage[margin=19mm]{geometry}
\usepackage{amsmath,amssymb,mathtools}
\usepackage[scr=rsfs,cal=euler]{mathalfa}
\usepackage{xfrac}
\usepackage[unicode,hidelinks,bookmarks=false]{hyperref}
\newcommand{\ie}{i.e.\ }
\newcommand{\cf}{cf.\ }
\newcommand{\resp}{respectively\ }
\newcommand{\Subsection}{\subsection}
\providecommand{\nobreakdash}{\nobreak\hbox{-}}
\setlength{\emergencystretch}{2em}
\multlinegap0pt
\newtheorem{theorem}{Theorem}
\newtheorem{definition}[theorem]{Definition}
\newtheorem{lemma}[theorem]{Lemma}
\newtheorem{corollary}[theorem]{CZrollary}
\newtheorem{conjecture}[theorem]{Conjecture}
\newtheorem{proposition}[theorem]{Proposition}
\newtheorem{remark}[theorem]{Remark}
\newcommand{\C}{\mathbb{C}}
\newcommand{\R}{\mathbb{R}}
\newcommand{\Z}{\mathbb{Z}}
\begin{document}
\title[New constructions relating Real and Complex Contact Structur]{New constructions relating Real and Complex Contact Structure}
\author{Ali M. Elgindi}
\date{}
\maketitle
\noindent\emph{Source and licence.} New constructions relating Real and Complex Contact Structure. Version arXiv:2607.01264v1 (CC BY 4.0). This material is distributed under the Creative Commons Attribution 4.0 International licence, http://creativecommons.org/licenses/by/4.0/, as recorded in the arXiv OAI licence metadata for this exact version and shown on the abstract page https://arxiv.org/abs/2607.01264v1. Full rights notice: licenses/arxiv-2607.01264-CC-BY-4.0.txt.\par\smallskip
\noindent\textit{Editorial scope.} Counted unit: the original \texttt{theorem} environment \texttt{-} at source lines 166--168 together with its complete proof at lines 170--219; the delivered document keeps source lines 165--220 so that every referenced label and every citation resolves inside it. Native editable source is the paper's own concatenated TeX; the AMS wrapper is editorial formatting only. No publisher class, upstream script or unrelated artwork is redistributed.
\section{Abundance of Complex Contact Structures}

\begin{theorem}[Uncountability of Complex Contact Structures]
The set of complex contact structures on $\mathbb{C}^3$ up to isotopy is uncountable.
\end{theorem}

\begin{proof}
We construct an uncountable family $\{\psi_\alpha\}_{\alpha\in I}$ of complex contact structures on $\mathbb{C}^3$ that are pairwise non-isotopic.

\medskip
Let $(S^3,\xi_{\mathrm{std}})$ be the standard contact 3-sphere. 
There exists an uncountable family of Legendrian knots $\{L_\alpha\}_{\alpha\in I}$ in $(S^3,\xi_{\mathrm{std}})$ 
that are smoothly isotopic but Legendrian non-isotopic. 
For instance, one may take a family of torus knots whose rotation numbers $\alpha\in\mathbb{R}/\mathbb{Z}$ 
are distinct. This is a standard result in Legendrian knot theory \cite{EtnyreHonda2001}.

\medskip
Choose an overtwisted contact structure $\xi_{\mathrm{ot}}$ on $S^3$ that is transverse to $\xi_{\mathrm{std}}$. 
Such a pair exists because overtwisted structures are flexible and can be chosen to be transverse to a given tight structure (see [4]). 
The space of contact structures on $S^3$ is large, and by the parametric $h$-principle for contact structures, 
we may assume that the pair $(\xi_{\mathrm{std}},\xi_{\mathrm{ot}})$ varies continuously with the parameter $\alpha$ 
when restricted to appropriate neighborhoods of the knots.

\medskip
For each Legendrian knot $L_\alpha$, we apply the corresponding explicit embedding as by our theorem \cite{Elgindi2025} to obtain an embedding
\[
F_\alpha: S^3 \hookrightarrow \mathbb{C}^3
\]
such that the complex tangents of $F_\alpha$ are exactly $L_\alpha$ and the holomorphic tangent planes along $L_\alpha$ 
are $\xi_{\mathrm{std}}|_{L_\alpha}$.

The pair $(\xi_{\mathrm{std}},\xi_{\mathrm{ot}})$ on the abstract $S^3$, together with the embedding $F_\alpha$, 
defines a \emph{formal complex contact structure} $\delta_\alpha$ along $F_\alpha(S^3)$ in $\mathbb{C}^3$. 
Concretely, the wedge construction $\delta_\alpha = \xi_{\mathrm{std}} \oplus J\xi_{\mathrm{ot}}$ gives an almost complex structure 
on a complex 2-plane bundle over $F_\alpha(S^3)$. considering this bundle together with the almost complex structure, will
constitute a formal solution of the complex contact relation in the sense of Forstnerič \cite{Forstneric2020}. 
The this formal construction varies continuously with $\alpha$ because the embeddings $F_\alpha$ and the contact structures can be chosen continuously.

\medskip
We now apply Forstnerič's parametric $h$-principle as given in Forstnerič \cite{Forstneric2020} which provided for a parametric $h$-principle of complex contact structures on Stein manifolds. 
Specifically, a continuous family of formal complex contact structures along a closed totally real submanifold $M\subset X$ 
(where $X$ is Stein) can be deformed into a family of genuine holomorphic complex contact structures on a Stein neighborhood of $M$, 
keeping the deformation fixed on a prescribed closed subset.

Here $X=\mathbb{C}^3$ (Stein) and $M = F_\alpha(S^3)$ (totally real except along $L_\alpha$). 
Applying the parametric $h$-principle to the family $\{\delta_\alpha\}$ yields a family $\{\psi_\alpha\}$ of genuine holomorphic complex contact structures 
defined on a Stein neighborhood of $F_\alpha(S^3)$ in $\mathbb{C}^3$. 
Because $\mathbb{C}^3$ is Stein and contractible, each $\psi_\alpha$ extends uniquely to all of $\mathbb{C}^3$.

\medskip
Suppose $\psi_\alpha$ and $\psi_\beta$ are isotopic as complex contact structures for $\alpha\neq\beta$. 
Then their formal data would be homotopic (by the $h$-principle), which would imply that the Legendrian knots $L_\alpha$ and $L_\beta$ 
are Legendrian isotopic in $(S^3,\xi_{\mathrm{std}})$. 
However, by construction $L_\alpha$ and $L_\beta$ have distinct rotation numbers and are therefore not Legendrian isotopic. 
This contradiction shows that $\{\psi_\alpha\}_{\alpha\in I}$ is an uncountable family of pairwise non-isotopic complex contact structures on $\mathbb{C}^3$.
\end{proof}


\begin{thebibliography}{99}
\bibitem[Elgindi2025]{Elgindi2025} A. M. Elgindi, \textit{Explicit holomorphic structures for embeddings of 3-manifolds into $\C^3$}, Methods Appl. Anal. \textbf{32} (2025), no. 1, 1–21.
\bibitem[EtnyreHonda2001]{EtnyreHonda2001} J. B. Etnyre and K. Honda, \textit{Knots and contact geometry I: Torus knots and the figure eight knot}, J. Symplectic Geom. \textbf{1} (2001), 63–120.
\bibitem[Forstneric2020]{Forstneric2020} F. Forstnerič, \emph{H-principle for complex contact structures on Stein manifolds}, Journal of Symplectic Geom. \textbf{18} (2020), no. 3, 733–767. arXiv:1810.12943.
\end{thebibliography}
\end{document}
